\begin{itemframe}{الگوریتم‌های تصادفی}
\itm
دیدیم که چگونه دانستن توزیع روی ورودی‌ها می‌تواند به ما در تحلیل رفتار یک الگوریتم در حالت میانگین کمک کند. اما اگر توزیع ورودی را ندانیم چه؟

\itm
آنگاه نمی‌توانیم الگوریتم را در حالت میانگین انجام تحلیل کنیم. پیش‌تر دیدم که در این حالت استفاده از یک الگوریتم تصادفی
\fn{randomized algorithms}
 می‌تواند مفید باشد.
\itm
برای برخی مسائل ،مانند مسئله‌ی استخدام، فرض برابر بودن احتمال همه‌ی جایگشت‌های ورودی مفید است.
در این حالت به جای آنکه یک توزیع برای ورودی فرض کنیم، یک توزیع به آن تحمیل می‌کنیم.
\end{itemframe}


\begin{itemframe}{الگوریتم‌های تصادفی}
\itm
دیدم که، پیش از اجرای الگوریتم استخدام، فهرست نامزدها را به صورت تصادفی جایگشت می‌دهیم تا تضمین کنیم که هر جایگشت به یک اندازه محتمل است. نشان دادیم با اینکار حدود
$\ln n$
بار یک دستیار جدید استخدام می‌کنیم.
\itm
با انجام این جایگشت تصادفی، اکنون این میانگین برای هر ورودی برقرار است، نه فقط برای ورودی‌هایی که از یک توزیع خاص گرفته شده‌اند.
\end{itemframe}


\begin{itemframe}{الگوریتم‌های تصادفی}
\itm
بیایید تفاوت میان تحلیل احتمالاتی و الگوریتم‌های تصادفی را بیشتر بررسی کنیم.
\itm
اگر الگوریتم استخدام به صورت قطعی پیاده شود، برای هر ورودی مشخص، تعداد دفعاتی که دستیار جدیدی استخدام می‌شود همیشه یکسان است.
چون این تعداد تنها به رتبه‌ها وابسته است، برای نمایش یک ورودی خاص کافی است رتبه‌ها را فهرست کنیم.
\itm
به طور مثال برای فهرست:
$$\langle 5,2,1,8,4,7,10,9,3,6 \rangle$$
، سه بار (رتبه‌های ۵، ۸، و ۱۰) دستیار جدید استخدام می‌شود.
\end{itemframe}


\begin{itemframe}{الگوریتم‌های تصادفی}
\itm
بنابراین هزینه الگوریتم، به ورودی وابسته است. در مقابل، الگوریتم تصادفی‌ای را در نظر بگیرید که ابتدا فهرست نامزدها را به صورت تصادفی جایگشت می‌دهد.
در این حالت برای یک ورودی خاص، نمی‌توانیم بگوییم چند بار مقدار بیشینه به‌روزرسانی می‌شود، چون این مقدار در هر بار اجرای الگوریتم فرق می‌کند.
\itm
برای بسیاری از الگوریتم‌های تصادفی، هیچ ورودی خاصی منجر به بدترین حالت نمی‌شود. حتی دشمن شما هم نمی‌تواند یک آرایه‌ی ورودی بد تولید کند، چون جایگشت تصادفی ترتیب ورودی را بی‌اثر می‌کند.
\itm
شبه کد مسئله استخدام که پیش‌تر بررسی کردیم را در نظر بگیرید. اگر در ابتدای آن، آرایه‌ی نامزدها را تصادفی جایگشت دهیم، الگوریتم
\texttt{RANDOMIZED-HIRE-ASSISTANT}
 به دست می آید.
\end{itemframe}


\begin{itemframe}{الگوریتم‌های تصادفی}
\itm
همانطور که قبلاً دیدیم، اگر ترتیب ورودی‌ها‌ تصادفی باشد هزینه‌ی استخدام در «حالت میانگین»
\texttt{HIRE-ASSISTANT}
برابر است با:
$$O(clnn)$$
\itm
در الگوریتم
\texttt{RANDOMIZED-HIRE-ASSISTANT}
بدون فرض توزیع خاصی بر روی ورودی هزینه استخدام «مورد انتظار» برابر است با
$$O(clnn)$$
\itm
اصطلاح، هزینه‌ی «حالت میانگین» را برای تحلیل احتمالاتی و هزینه‌ی «مورد انتظار» را برای الگوریتم‌های تصادفی به کار می‌بریم. وقتی می‌گوییم حالت میانگین، یعنی تحلیل احتمالاتی انجام داده‌ایم و توزیعی برای ورودی فرض کرده‌ایم اما در مورد الگوریتم‌های تصادفی حدس توزیع ورودی وجود ندارد پس بهتر است این دو کار را با اسامی متفاوتی نام‌گذاری کنیم.
\end{itemframe}


\begin{itemframe-s}{الگوریتم‌های تصادفی}{جایگشت تصادفی آرایه‌ها}
\itm
بسیاری از الگوریتم‌های تصادفی، ورودی را به صورت تصادفی جایگشت می‌دهند.
هدف تولید یک «جایگشت تصادفی یکنواخت»
\fn{uniform random permutation}
است، یعنی هر جایگشت به یک اندازه محتمل باشد.
\itm
چون تعداد جایگشت‌ها برابر با $n!$ است، می‌خواهیم احتمال تولید هر جایگشت برابر با
$\frac{1}{n!}$
باشد.
\end{itemframe-s}


\begin{itemframe-s}{الگوریتم‌های تصادفی}{جایگشت تصادفی آرایه‌ها}
\itm
الگوریتم
\texttt{RANDOMLY-PERMUTE}
 آرایه
$A[1 \dots n]$
را در زمان
$\Theta(n)$
جایگشت می‌دهد:
\begin{algorithm}[H]\alglr
\caption{RANDOMLY-PERMUTE}
\begin{algorithmic}[1]
\Procedure{RANDOMLY-PERMUTE}{$A, n$}
  \For{$i \gets 1$ to $n$}
    \State swap A[i] with A[RANDOM(i, n)]
  \EndFor
\EndProcedure
\end{algorithmic}
\end{algorithm}

\itm
این الگوریتم یک جایگشت تصادفی یکنواخت تولید می‌کند.
\end{itemframe-s}